\section{失败分析、结果解释与适用局限}
\label{cr:discussion}
\subsection{工程完成与科学有效性的区别}
预充失败表明，完成归档、求解器正常结束和读出分类为正确是三个不同判断。
一个运行即使到达$4\ns$，仍可能未满足预先规定的电学条件。
旧$100\,\Omega$事件的读出裕量远离错误边界，但其预充状态超差，
因此仍不能作为有效主响应。R50修订改变了外部电路，而不是更改误差容限；
它需要独立版本、重新初始化和新的匹配参考，不能称为单纯数值收敛改进。

历史过程还出现输入名称导致的解析错误、数据收集中断和工程重试。
处理原则是保留原输入、日志和失败身份，先确认远端实际状态，再恢复收集或
进行获准的唯一重试。成功重试只补足原逻辑槽位，不增加粒子事件数。
这类证据适合写入附录以说明流程可靠性；正文应着重解释它对数据纳入和结论的影响，
避免用大量命令输出取代研究论证。

\subsection{电荷变化、参考响应与分类的关系}
存储节点在粒子作用后降低，不意味着存储0和存储1的读出裕量都会变差。
式\eqref{cr:metrics}对位值引入符号后，同方向的电压扰动可能削弱存储1，
却增大存储0的裕量。源注入总电荷也不等于存储电容的净变化：
载流子可被不同端口收集、复合，并通过寄生电容与外电路交换电荷。
因此，不能以$Q_{\mathrm{dep}}/C$直接替代真实瞬态求解，也不应把固定电容下
电压和电荷的等价换算当作两项独立验证。

D1进一步说明，参考支路决定了哪些共模响应会留在差分中。
理想固定参考不会跟随访问管的读出耦合，而匹配中性支路会响应同样的字线操作。
但匹配并非完美消除所有状态依赖：信号与参考的初态不同，其沟道和结区演化
也可能不同。R50采用物理双支路一致性检查，是为了确认离线差分的实现，
并未因此完成真实感放失调、噪声和锁存恢复的校准。

\subsection{预测效果的含义与潜在偏差}
训练交叉验证明显优于均值基准，支持继续检验冻结代理模型；
它尚不能证明模型在$0.45\um$或新能量条件下仍有相同精度。
特征包含已知Geant4沉积量，使模型获得粒子与路径差异的重要信息，
也限定了预测的用途：它替代的是后续TCAD响应近似，不是完整输运过程。
当前仅有两个粒子类型、一个能量和两个方向，不能据此讨论连续空间辐射谱。

分组交叉验证减少了同工况重复跨折导致的乐观评价，但样本仍来自共同的简化模型。
未校准几何、理想电容、受控初态和统一材料假设可能共同形成稳定的系统偏差，
这种偏差不会因训练误差变小而被发现。
后续测试应同时报告逐点残差、分层误差及基准比较；若外推误差明显增大，
应分析特征与物理状态的适用边界，而非通过回看测试集调整模型后仍称原测试有效。

\subsection{仍需完成的技术与课程工作}
当前优先完成已冻结测试响应的真实求解和评分，再执行既定4/6 nm氧化层变体、
主矩阵最近边界点的匹配细化，以及R50原生场图的实际视觉验收。
氧化层扩展只讨论固定硅入口事件下的TCAD厚度敏感性；
Geant4仍为硅/真空输运，不能把它写成完整材料栈工艺变化。
电磁option4未包含核反应物理，质子结果不能解释为已验证的核反冲SEE。

若进一步面向器件研究，需要获得具有出处的DRAM结构和电学参考，
校准访问管、写入保持、感放与恢复，并按应用目标扩展粒子能量、位置与统计规模。
若面向课程验收，还需学生亲自解释模型、运行限定演示、完成真实分工与教师审核。
AI生成脚本和报告可以支撑这些活动，但不能替代参与人数、学生理解或考核事实。
现阶段可交付的是已验证技术产物及其明确边界，最终课程结论应由完整技术与教学
证据共同支持。
